% !TEX root = ../main.tex

\section{Silnik Otto - sprawność w punkcie maksymalnej mocy}
Podczas cyklu Otto ciepło jest wymieniane jedynie w przemianach izochorycznych. Szybkość przepływu ciepła jest proporcjonalna do różnicy temperatur między gazem roboczym a zbiornikiem ciepła zgodnie z prawem Fouriera oraz z warunkami brzegowymi \cite{endo_otto}:
\begin{align}
	\odv{T}{t} &= -\alpha_h \pab{T - T_h}, &
	T\pab{t=0} &= T_2, &
	T\pab{t=\tau_h} &= T_3, \\
	\odv{T}{t} &= -\alpha_c \pab{T - T_c}, &
	T\pab{t=0} &= T_4, &
	T\pab{t=\tau_c} &= T_1,
\end{align}
gdzie \(T_h\) oraz \(T_c\) to odpowiednio temperatury zbiornika ciepłego i zimnego, które nie są osiągane przez układ; temperatury \(T_1\), \(T_2\), \(T_3\) oraz \(T_4\) to temperatury w odpowiednich punktach wykresu \ref{fig:pv_silnik_otto}; wielkości \(\alpha_h\) oraz \(\alpha_c\) to współczynniki opisujące przenikalność cieplną ścianek układu. Nie są to te same współczynniki, które pojawiają się w analizie cyklu Carnota. Nie są one ze sobą bezpośrednio porównywane i nie prowadzi to do konfliktu oznaczeń. Rozwiązując równania różniczkowe z zadanymi warunkami brzegowymi otrzymujemy:
\begin{equation}
\begin{aligned}
	T_3 - T_h &= \pab{T_2 - T_h} e^{-\alpha_h \tau_h}, &
	T_1 - T_c &= \pab{T_4 - T_c} e^{-\alpha_c \tau_c}, \\
	T_2 - T_h &= \pab{T_3 - T_h} e^{\alpha_h \tau_h}, &
	T_4 - T_c &= \pab{T_1 - T_c} e^{\alpha_c \tau_c}.
\end{aligned}
\end{equation}
Możemy je przepisać w postaci:
\begin{equation}\label{eq:otto_endo_T_roznice}
\begin{aligned}
	T_3 - T_2 e^{-\alpha_h \tau_h} &= T_h \pab{1 - e^{-\alpha_h \tau_h}}, &
	T_1 - T_4 e^{-\alpha_c \tau_c} &= T_c \pab{1 - e^{-\alpha_c \tau_c}}, \\
	T_3 e^{\alpha_h \tau_h} - T_2 &= T_h \pab{e^{\alpha_h \tau_h} - 1}, &
	T_1 e^{\alpha_c \tau_c} - T_4 &= T_c \pab{e^{\alpha_c \tau_c} - 1}.
\end{aligned}
\end{equation}
Równanie \eqref{eq:otto_stosunek_T} wyznaczyliśmy z równań przemian adiabatycznych i są prawdziwe również w tym przypadku i zapisujemy je w formie:
\begin{equation}\label{eq:otto_endo_stosunek_T}
	\frac{T_1}{T_2} = \frac{T_4}{T_3} = \inv{r^{\varkappa-1}} = \kappa,
\end{equation}
gdzie \(\kappa\) została wprowadzona w celu uproszczenia wyrażeń.

\begin{figure}[htp]
	\centering
	\tikzsetnextfilename{pv_silnik_endo_otto_70}
\begin{tikzpicture}
	\begin{axis}[width=0.7\textwidth,
		xlabel=$V$,
		ylabel=$p$,
		axis x line=bottom,
		axis y line=middle,
		xmin=0,
		xmax=1.2,
		ymin=0,
		ymax=1.2,
		ticks=none,
		every axis x label/.style={at={(current axis.right of origin)},anchor=north east},
		every axis y label/.style={at={(current axis.above origin)},anchor=north east},
		declare function={
			r = 2;
			tau = 2;
			f = 3;
			kappa = 1 + 2/f;
			V1 = 1/r;
			V2 = 1;
			p1 = 1/r^kappa * 1/tau;
			p2 = 1/tau;
			p3 = 1;
			p4 = 1/r^kappa;
		}]
		\addplot[mark=*] coordinates {(V2,p1)} node[anchor=north west]{$1$};
		\addplot[mark=*] coordinates {(V1,p2)} node[anchor=north east]{$2$};
		\addplot[mark=*] coordinates {(V1,p3)} node[anchor=south east]{$3$};
		\addplot[mark=*] coordinates {(V2,p4)} node[anchor=south west]{$4$};
		\addplot[domain=V1:V2] {1/r^kappa * 1/tau / x^kappa} node [pos=0.5, below, sloped] {$Q=0$};
		\addplot[mark=none] coordinates {(V1, p2) (V1, p3)};
		\addplot[domain=V1:V2] {1/r^kappa / x^kappa}         node [pos=0.5, above, sloped] {$Q=0$};
		\addplot[mark=none] coordinates {(V2, p4) (V2, p1)};
		\addplot[mark=none, dashed] coordinates {(V2, p1) (V2, 0)};
		\addplot[mark=none, dashed] coordinates {(V1, p2) (V1, 0)};
		\draw[line width=0.3mm, -{Stealth[length=5mm, open]}] (0.4,0.7) node[anchor=east]{$Q_{2-3}$} -- (0.5,0.7);
		\draw[line width=0.3mm, -{Stealth[length=5mm, open]}] (1,0.25) -- (1.1,0.25) node[anchor=north]{$Q_{4-1}$};
		
		\addplot[domain=0.9*V1:1.25*V1, dashed] {1.1*1/r * 1 / x}             node [pos=0.7, above, sloped] {$T_h$};
		\addplot[domain=0.9*V2:1.1*V2, dashed] {0.6*1/r^kappa * 1/tau * 1 / x}             node [pos=0.3, below, sloped] {$T_c$};
	\end{axis}
\end{tikzpicture}

	\caption{Wykres cyklu Otto w warunkach endoodwracalnych}
\end{figure}

Ze wzorów \eqref{eq:otto_endo_T_roznice} oraz \eqref{eq:otto_endo_stosunek_T} możemy wyprowadzić wzory na temperaturę zależne od \(\tau_h\), \(\tau_c\) oraz \(\kappa\):
\begin{equation}\label{eq:endo_otto_T}
	\begin{split}
		T_1 &= \frac{
			T_c e^{\alpha_h \tau_h} \pab{e^{\alpha_c \tau_c} - 1} + \kappa T_h \pab{e^{\alpha_h \tau_h} - 1}
		}{
			e^{\alpha_h \tau_h + \alpha_c \tau_c} - 1
		}, \\
		T_2 &= \frac{
			T_c e^{\alpha_h \tau_h} \pab{e^{\alpha_c \tau_c} - 1} + \kappa T_h \pab{e^{\alpha_h \tau_h} - 1}
		}{
			\kappa \pab{e^{\alpha_h \tau_h + \alpha_c \tau_c} - 1}
		}, \\
		T_3 &= \frac{
			T_c \pab{e^{\alpha_c \tau_c} - 1} + \kappa T_h e^{\alpha_c \tau_c} \pab{e^{\alpha_h \tau_h} - 1}
		}{
			\kappa \pab{e^{\alpha_h \tau_h + \alpha_c \tau_c} - 1}
		}, \\
		T_4 &= \frac{
			T_c \pab{e^{\alpha_c \tau_c} - 1} + \kappa T_h e^{\alpha_c \tau_c} \pab{e^{\alpha_h \tau_h} - 1}
		}{
			e^{\alpha_h \tau_h + \alpha_c \tau_c} - 1
		}.
	\end{split}
\end{equation}

Sprawność cyklu Otto w punkcie maksymalnej mocy jest identyczna z przypadkiem równowagowym i zależne jedynie od \(\kappa\):
\begin{equation}
	\eta = 1 - \inv{r^{\varkappa-1}} = 1 - \kappa.
\end{equation}
Ciepła dostarczone oraz odprowadzone są równe:
\begin{gather}
	Q_{2-3} = nc_V \pab{T_3 - T_2} =
	nc_V \frac{
	\pab{\kappa T_h - T_c}
	\pab{e^{\alpha_h \tau_h} - 1}
	\pab{e^{\alpha_c \tau_c} - 1}
	}{
	\kappa\pab{e^{\alpha_h \tau_h + \alpha_c \tau_c} - 1}
	}, \\
	Q_{4-1} = nc_V \pab{T_1 - T_4} =
	- nc_V \frac{
	\pab{\kappa T_h - T_c}
	\pab{e^{\alpha_h \tau_h} - 1}
	\pab{e^{\alpha_c \tau_c} - 1}
	}{
	e^{\alpha_h \tau_h + \alpha_c \tau_c} - 1
	}.
\end{gather}
Praca wykonana w jednym cyklu według wzoru \eqref{eq:praca_w_silniku}:
\begin{multline}
	-W = Q_{2-3} + Q_{4-1} =
	\frac{nc_V\pab{1-\kappa}\pab{\kappa T_h - T_c}}{\kappa}
	\frac{
		\pab{e^{\alpha_h \tau_h} - 1}
		\pab{e^{\alpha_c \tau_c} - 1}
	}{
		e^{\alpha_h \tau_h + \alpha_c \tau_c} - 1
	} = \\
	\frac{2nc_V\pab{1-\kappa}\pab{\kappa T_h - T_c}}{\kappa}
	\sinh\pab{\frac{\alpha_h \tau_h}{2}}
	\sinh\pab{\frac{\alpha_c \tau_c}{2}}
	\csch\pab{\frac{\alpha_h \tau_h + \alpha_c \tau_c}{2}}.
\end{multline}
Moc jest stosunkiem pracy w jednym cyklu do długości tego cyklu \(\tau = \zeta\pab{\tau_h + \tau_c}\) i w rozpatrywanym przypadku dana jest wzorem \cite{endo_otto}:
\begin{equation}
	P = \frac{nc_V}{\zeta\pab{\tau_h + \tau_c}}
	\frac{\pab{1-\kappa}\pab{\kappa T_h - T_c}}{\kappa}
	\sinh\pab{\frac{\alpha_h \tau_h}{2}}
	\sinh\pab{\frac{\alpha_c \tau_c}{2}}
	\csch\pab{\frac{\alpha_h \tau_h + \alpha_c \tau_c}{2}}.
\end{equation}

Wyrażenie na sprawność silnika Otto zależy jedynie od \(\kappa\), natomiast na moc można rozseparować na czynnik zależny od \(\kappa\) i od pozostałych zmiennych. Aby znaleźć sprawność w punkcie maksymalnej mocy wystarczy znaleźć maksimum funkcji zależnej od \(\kappa\):
\begin{gather*}
	f\pab{\kappa} = \frac{1-\kappa}{\kappa}\pab{\kappa T_h - T_c}, \\
	\odv[delims-eval=.|]{f}{\kappa}_{\kappa_{max}} = \inv{\kappa_{max}^2}T_c - T_h = 0, \\
	\kappa_{max} = \sqrt{\frac{T_c}{T_h}}.
\end{gather*}
Dla tej wartości parametru \(\kappa\) sprawność jest równa:
\begin{equation}\label{eq:endo_otto_klas_eta}
	\eta\pab{\kappa_{max}} = 1 - \sqrt{\frac{T_c}{T_h}}.
\end{equation}
Jest to wynik identyczny do sprawności cyklu Carnota w punkcie maksymalnej mocy \eqref{eq:endo_carnot_eta}.
